\documentclass[11pt]{article}
\usepackage[margin=1in]{geometry}
\usepackage{amsmath,amssymb}
\usepackage{enumitem}
\usepackage[hidelinks]{hyperref}

\title{APM348 Project Outline}
\author{Neel Sarkar \and Piper Scott}
\date{February 2026}

\begin{document}
\maketitle

\noindent
\textbf{Title:} Content-Centric Modeling of Rage-Bait Lifecycles on Social Media with Toxicity Feedback\\
\textbf{Course:} APM348 (Winter 2026), University of Toronto

\section*{1. Background and Motivation}
This project studies how platform ranking can drive a tradeoff between short-run engagement and long-run platform health. Prior evidence shows engagement-based algorithms can amplify divisive content \cite{huszar2022,milli2025}, while toxic content can affect user behavior \cite{beknazar2025}. Moral-emotional content also spreads more readily \cite{brady2017}. We model this using a compartmental ODE system where \textbf{content items} (not users) are the primary dynamical units.

\section*{2. Previous Mathematical Models}
SIR-style models and rumor models \cite{daleyKendall1964} give threshold and stability tools for diffusion, and recent work applies these ideas to social-media spread \cite{liShao2024}. For example, the SVFR model models users' decisions, and so focus on local networks. Toxicity-focused compartment models also exist \cite{maleki2022,addai2025}. The main gap is that most models compartmentalize \emph{users}; our model compartmentalizes \emph{content lifecycle states} and couples them to toxicity and user attrition, describing a more global scale.

\section*{3. Proposed Model and Planned Contributions}
State variables are $I,V,F,S,\tau,U$ for ignored, viral, fatigued, moderated/suppressed content, toxicity, and active users. Core transitions are:
\[
\dot{I}=\rho U-\beta_{\mathrm{eff}}IV-(\delta+\mu_c)I,\quad
\dot{V}=\beta_{\mathrm{eff}}IV-\gamma_{\mathrm{eff}}V-\mu_cV,\quad
\dot{F}=\gamma_{\mathrm{eff}}V-\mu_cF,\quad
\dot{S}=\delta I,
\]
\[
\dot{\tau}=\phi V-\psi\tau,\qquad
\dot{U}=\nu-\lambda_u(1+w\tau)U,
\]
with
\[
\beta_{\mathrm{eff}}=\alpha\beta_0(1+\kappa\tau),\qquad
\gamma_{\mathrm{eff}}=\gamma_0(1+\eta\tau).
\]
Here $\rho$ is content injection per user, $\mu_c$ natural content decay, and $\delta$ moderation/suppression. Main contribution: a content-centric model that links amplification, toxicity, and user-base dynamics in one system.
For simplicity, direct moderation is modeled only from $I$ via $\delta I$; an extension can include moderation from $V$ or $F$.
Natural decay $\mu_c$ represents loss of visibility/attention (content leaving the system), while $S$ tracks explicit moderation removals.
For the empirical Higgs fit, we use a reduced IVF calibration model for the retweet cascade with a decaying exogenous seed term $\lambda(t)=\lambda_0 e^{-d_\lambda t}$. This improves the data fit for the observed tail without changing the main IVFS policy model written above.

\section*{4. Planned Analysis and Simulations}
\begin{itemize}[leftmargin=1.3em]
\item \textbf{Calibration data:} SNAP Higgs activity stream \cite{snap}. Verified counts: 563,069 total events, 354,930 retweets, peak at hour 78 with 30,602 RTs.
\item \textbf{Parameter fitting:} fit a reduced IVF calibration model to a 100-hour normalized retweet window (scale-invariant) using a bounded Powell search with warm starts. The fitted calibration parameters are $\beta_0,\gamma_0,\lambda_0,d_\lambda,V_0$. The full IVFS scenario parameters are then held fixed when we run the policy model. At the moment, the toxicity couplings $\phi,\psi$ are not jointly calibrated from Higgs; Jigsaw is used only as an external reference check, so the toxicity side should be interpreted as scenario-driven rather than fully data-calibrated.
\item \textbf{Analytical work:} derive equilibria, a disease-free threshold analogue $\mathcal{R}_0$, and local Jacobian stability conditions.
\item \textbf{Numerical scenarios:} engagement-first ($\alpha=0.9$), moderate ($\alpha=0.5$), health-first ($\alpha=0.2$), with long horizons (e.g., $t=1000$) for steady values of $V^*,\tau^*,U^*$ and $E=\alpha V^*U^*$. Under the current calibration, the threshold occurs at $\alpha(\mathcal{R}_0=1)\approx 0.316$, so the health-first case is subcritical while the moderate and engagement-first cases remain above threshold.
\item \textbf{Continuation sweep:} compute long-run $\tau^*(\alpha)$ for $\alpha\in[0.05,1.0]$ and report as continuation evidence (not a formal bifurcation claim unless stability checks support it).
\item \textbf{Sensitivity:} vary $\kappa,\eta,\phi,\psi,\lambda_u$ to quantify robustness, with particular attention to $\phi$ because toxicity is still scenario-driven. A dedicated $\phi$-sweep now serves as a core robustness check: with $\psi$ fixed, the current model value $\phi=0.056$ is compared against the larger external Jigsaw-implied value $\phi\approx 0.146$. The main qualitative ordering remains stable across this range: engagement-first yields the highest long-run toxicity, moderate remains in the middle, and the health-first case stays suppressed because it is already below the calibrated $\mathcal{R}_0=1$ threshold.
\item \textbf{Uncertainty quantification:} parametric Poisson bootstrap ($B=500$) gives 95\% CIs for $\beta_0\in[0.9983,1.0021]$, $\gamma_0\in[0.2962,0.3044]$, and $\alpha(\mathcal{R}_0\!=\!1)\in[0.3126,0.3205]$. Local profile-likelihood curves around the optimum show clear minima for both $\beta_0$ and $\gamma_0$, supporting practical identifiability under the current calibration model.
\end{itemize}

\section*{5. Expected Results and Conclusion Plan}
Current results indicate that higher amplification improves short-run virality but worsens long-run toxicity and user retention. They also suggest a meaningful threshold effect: once $\alpha$ drops below the calibrated $\mathcal{R}_0=1$ boundary, viral activity is no longer self-sustaining. The new $\phi$-robustness check does not remove the toxicity-calibration limitation, but it does show that the main policy ranking is not an artifact of one arbitrary $\phi$ choice. The conclusion will summarize equilibrium and stability findings, compare policy scenarios, discuss implications for platform moderation and ranking policy, and state limitations clearly.

\paragraph{Limitations.}
Homogeneous mixing, deterministic dynamics, and single-index toxicity are the main structural simplifications.
The fitted $\beta_0,\gamma_0$ and calibration-side seed parameters are estimated from a single Higgs cascade, so the empirical spread fit is stronger than the toxicity fit. In particular, toxicity is still scenario-driven rather than jointly calibrated from the same dataset, and the Jigsaw comparison should be treated as an external reference rather than a direct platform match. The current $\phi$-sweep should therefore be read as a robustness argument, not as a substitute for true same-dataset toxicity calibration. Even so, the bootstrap intervals are narrow and the local profile curves show well-defined minima for both $\beta_0$ and $\gamma_0$, lending practical identifiability support for the spread component.
Extensions to network-aware and stochastic models are natural next steps.

\begingroup
\footnotesize
\begin{thebibliography}{9}

\bibitem{snap}
Stanford Network Analysis Project.
Higgs Twitter Dataset.
\url{https://snap.stanford.edu/data/higgs-twitter.html}, 2013.

\bibitem{daleyKendall1964}
D.~J.~Daley and D.~G.~Kendall.
Epidemics and rumours.
\textit{Nature}, 204:1118, 1964.

\bibitem{brady2017}
W.~J.~Brady, J.~A.~Wills, J.~T.~Jost, J.~A.~Tucker, and J.~J.~Van~Bavel.
Emotion shapes the diffusion of moralized content in social networks.
\textit{Proc.\ Natl.\ Acad.\ Sci.\ USA}, 114(28):7313--7318, 2017.

\bibitem{huszar2022}
F.~Husz\'{a}r, S.~I.~Ktena, C.~O'Brien, L.~Belli, A.~Schlaikjer, and M.~Edizel.
Algorithmic amplification of politics on Twitter.
\textit{Proc.\ Natl.\ Acad.\ Sci.\ USA}, 119(1):e2025334119, 2022.

\bibitem{milli2025}
S.~Milli, M.~Carroll, Y.~Wang, S.~Pandey, S.~Zhao, and A.~Dragan.
Engagement, user satisfaction, and the amplification of divisive content on social media.
\textit{PNAS Nexus}, 4(1):pgae594, 2025.

\bibitem{beknazar2025}
G.~Beknazar-Yuzbashev, R.~Jim\'{e}nez~Dur\'{a}n, J.~McCrosky, and M.~Stalinski.
Toxic content and user engagement on social media: Evidence from a field experiment.
CESifo Working Paper 11644, 2025.

\bibitem{liShao2024}
Y.~Li and L.~Shao.
Using an epidemiological model to explore the interplay between sharing and advertising in viral videos.
\textit{Scientific Reports}, 14:11351, 2024.

\bibitem{maleki2022}
M.~Maleki, V.~Mead, and N.~Agarwal.
Applying an epidemiological model to evaluate the propagation of toxicity on Twitter.
In \textit{Proc.\ HICSS-55}, 2022.

\bibitem{addai2025}
E.~Addai, N.~Agarwal, and N.~Yousefi.
Modeling and simulation of interventions' effect on the spread of toxicity in social media.
\textit{Online Social Networks and Media}, 46:100309, 2025.

\bibitem{jing2018}
Y.~Jing, H.~Wang, and Z.~Zhang.
Improved SIR advertising spreading model and its effectiveness in social network.
\textit{Procedia Computer Science}, 129:215--218, 2018.

\end{thebibliography}
\endgroup

\end{document}
